\documentclass[12pt]{article}
%fixes to the latex2e kernel
%\usepackage{fixltx2e} %this is not needed after 2015
\usepackage{fix-cm}
\usepackage{etex}

%fix double floats numbering and positioning
\usepackage{dblfloatfix}

%checks for obsolete packages
\usepackage{nag}


%%% Local Variables: 
%%% mode: latex
%%% End: 

%Note: pagination needs to be loaded after graphics, because mdframed
%needs to be loaded after xcolor to keeep the our options for the latter
%colors
\makeatletter
\@ifpackageloaded{xcolor}{}{%
\usepackage[table,x11names,dvipsnames,svgnames]{xcolor}%
}
\makeatother

%colors in table
\usepackage{colortbl}

%pdf
\usepackage{graphicx}
\usepackage{wrapfig}

% Lyft colors (see https://design.lyft.com/re-approaching-color-9e604ba22c88)
\input{preamble/graphicsColors}

%%% Local Variables: 
%%% mode: latex
%%% End: 

\usepackage{cite}

%advanced typesetting
\usepackage{microtype}

% US hyphenation
\usepackage[american]{babel}


% extensions for tables
\usepackage{array}
\usepackage{multirow}
\usepackage{booktabs}
\usepackage{makecell} %introduces \thead and \makecell

%compact paragraph title
\newcommand{\myparagraph}[1]{\textbf{\emph{#1}}.}
%\newcommand{\subparagraph}[1]{\emph{#1}.}

%provide options for changing spacing in enumeration environments
\ifcsname labelindent\endcsname
\let\labelindent\relax
\fi
\usepackage[inline]{enumitem}

%provides subfloats (subcaption replaces subfig and subfigure, but
%might not be compatible with some classes)
\usepackage{subfig}

%Inspired by CVPR template
\providecommand{\ie}{i.e\onedot}
\newcommand{\eg}{e.g\onedot}
\providecommand{\etal}{et al\onedot}

%set more relaxed constraints on the floats
\setcounter{topnumber}{2}
\setcounter{bottomnumber}{2}
\setcounter{totalnumber}{4}
\renewcommand{\topfraction}{0.85}
\renewcommand{\bottomfraction}{0.85}
\renewcommand{\textfraction}{0.15}
\renewcommand{\floatpagefraction}{0.7}

%Make an enumeration with a letter+progressive number
\newenvironment{lenumerate}[2][]
{\begin{enumerate}[label=(#2\arabic*),leftmargin=0.2in,itemindent=0.15in,#1]}
{\end{enumerate}}

%Make an letter+progressive number description list
\newenvironment{ldescription}[2][]
{\begin{lenumerate}{#2}%
\let\bitem\item%
\renewcommand{\item}[1][]{\bitem\textsl{##1}:~}}
{\end{lenumerate}}


 %The following sets the labeling for inline enumerations

\setlist*[enumerate,1]{label={\itshape\arabic*)}}

%Define macro to make paragraph headings always end with a full stop
\makeatletter
\newcommand{\paragraphswithstop}{%
\let\copyparagraph\paragraph%
\renewcommand\paragraph[1]{\copyparagraph{##1.}}%
}
\makeatother

%Package to frame text in boxes
\usepackage[framemethod=tikz]{mdframed}

% References by text instead of section or page number
\makeatletter
\def\namedlabel#1#2{\begingroup
  #2%
  \def\@currentlabel{#2}%
  \phantomsection\label{#1}\endgroup
}
\makeatother

%Same as above, but do not insert the label text
\makeatletter
\def\namedlabelphantom#1#2{\begingroup
  \def\@currentlabel{#2}%
  \phantomsection\label{#1}\endgroup
}
\makeatother

%Force exact termination of a paragraph with justification
\newcommand{\parunskip}{\bgroup\unskip\parfillskip=0pt \par\egroup}

%%% Local Variables:
%%% mode: latex
%%% End:

%
% AMS typesetting
\makeatletter
\@ifpackageloaded{amsmath}{}{%
\usepackage{amsmath}%
}
\makeatother
\usepackage{mathtools}
\usepackage{amssymb,amsfonts}
\makeatletter
\@ifundefined{proof}{%
\usepackage{amsthm}%
}{}
\makeatother
% Generic symbols (e.g., \degree)
\usepackage{gensymb}
% Prescripts, squashed summations, and others
\usepackage{mathtools}
% Bold greek letters
\usepackage{bm}

%Common math statements environments
\makeatletter
\@ifundefined{theorem}{
\newtheorem{theorem}{Theorem}
\newtheorem{corollary}{Corollary}
\newtheorem{proposition}{Proposition}
\newtheorem{lemma}{Lemma}
\newtheorem{fact}{Fact}
\newtheorem{problem}{Problem}
\newtheorem{assumption}{Assumption}
\newtheorem{property}{Property}
\newtheorem{remark}{Remark}
\newtheorem{example}{Example}
}{}
  %SIAM article classes give their own Definition and Remark environments
\@ifundefined{definition}{
  \newtheorem{definition}{Definition}
}
% \@ifundefined{remark}{
%   \newtheorem{remark}{Remark}
% }
\makeatother
\newcommand{\rmss}[1]{_{\textrm{#1}}}

%%% Local Variables:
%%% mode: latex
%%% End:

%Spaces
\newcommand{\real}[1]{\mathbb{R}^{#1}{}}
\newcommand{\complex}[1]{\mathbb{C}^{#1}{}}
\newcommand{\naturals}[1]{\mathbb{N}^{#1}{}}
\newcommand{\integers}[1]{\mathbb{Z}^{#1}{}}
\newcommand{\sphere}[1]{{\mathbb{S}^{#1}}{}}
\newcommand{\stiefel}{\mathrm{St}}
\newcommand{\grassmann}{\mathrm{Grass}}
\newcommand{\GL}{\mathbb{GL}}

%short-hand for matrices
\newcommand{\bmat}[1]{\begin{bmatrix}#1\end{bmatrix}}
\newcommand{\Bmat}[1]{\begin{Bmatrix}#1\end{Bmatrix}}
\newcommand{\pmat}[1]{\begin{pmatrix}#1\end{pmatrix}}
\newcommand{\smallbmat}[1]{\left[\begin{smallmatrix}#1\end{smallmatrix}\right]}

%supertscript operators
\newcommand{\transpose}{^\mathrm{T}}
\newcommand{\hermitian}{^\mathrm{H}}
\newcommand{\inverse}{^{-1}}
\newcommand{\pinverse}{^\dagger}
\newcommand{\orth}{^{\bot}}
\newcommand{\apstar}{^{\ast}}

%parentheses-based operators
\newcommand{\cross}[1]{[#1]_{\times}\!}
\newcommand{\crossinv}[1]{[#1]_{\times}^{\textrm{inv}}\!}

%equality
\newcommand{\defeq}{\doteq}


%Norms, absolute values, and inner products
\DeclarePairedDelimiter{\abs}{\lvert}{\rvert}
\DeclarePairedDelimiter{\ceil}{\lceil}{\rceil}
\DeclarePairedDelimiter{\norm}{\lVert}{\rVert}
\newcommand{\frob}[1]{\norm{#1}_F}
\newcommand{\bigfrob}[1]{\bignorm{#1}_F}
\newcommand{\metric}[3][]{g_{#1}\left(#2, #3\right)}
\newcommand{\metrica}[3][]{\langle #2, #3\rangle_{#1}}

%Derivatives
\newcommand{\de}{\mathrm{d}}
\newcommand{\dert}[1][]{\frac{\de #1}{\de t}}
\newcommand{\ddert}[1][]{\frac{\de^2 #1}{\de t^2}}
%Vector
\newcommand{\vct}[1]{\mathbf{#1}}
\newcommand{\pder}[2][]{\frac{\partial #1}{\partial #2}}
%named operators
\DeclareMathOperator*{\Land}{\bigwedge}
\DeclareMathOperator*{\Lor}{\bigvee}
\DeclareMathOperator{\rank}{rank}
\DeclareMathOperator{\diag}{diag}
\DeclareMathOperator{\blkdiag}{blkdiag}
\DeclareMathOperator{\symm}{sym}
\DeclareMathOperator{\asym}{skew}
\DeclareMathOperator*{\argmin}{\arg\!\min}
\DeclareMathOperator*{\argmax}{\arg\!\max}
\DeclareMathOperator{\softmin}{softmin}
\DeclareMathOperator{\softmax}{softmax}
\DeclareMathOperator{\trace}{tr}
\DeclareMathOperator{\vecspan}{span}
\DeclareMathOperator{\vecnull}{null}
\DeclareMathOperator{\vecnullity}{nullity}
\DeclareMathOperator{\vecop}{vec}
\DeclareMathOperator{\stack}{stack}
\DeclareMathOperator{\hstack}{hstack}
\DeclareMathOperator{\vstack}{vstack}
\DeclareMathOperator{\sign}{sign}
\DeclareMathOperator{\sinc}{sinc}

\newcommand{\compose}{\circ}

\newcommand{\divergence}{\nabla\cdot}
\DeclareMathOperator{\grad}{{grad}}
\DeclareMathOperator{\D}{D\!}
\DeclareMathOperator{\Dgrad}{{Dgrad}}
\DeclareMathOperator{\hess}{{Hess}}


\DeclareMathOperator{\expect}{\mathbb{E}}
\DeclareMathOperator{\probability}{\mathbb{P}}

\DeclareMathOperator{\expm}{expm}
\DeclareMathOperator{\logm}{logm}
\DeclareMathOperator{\Log}{Log}
\DeclareMathOperator{\LogNorm}{\frac{Log}{\norm{Log}}}
\DeclareMathOperator{\DLog}{DLog}
\DeclareMathOperator{\DLogNorm}{D\frac{Log}{\norm{Log}}}

\DeclareMathOperator{\cvxhull}{co}
\DeclareMathOperator{\relu}{ReLU}
\DeclareMathOperator{\proj}{proj}
\DeclareMathOperator{\dist}{dist}

\DeclareMathOperator*{\bigO}{O}

%\DeclareMathOperator*{\Pr}
\newcommand{\intersect}{\cap}
\newcommand{\union}{\cup}
\DeclareMathOperator*{\Or}{\bigvee}

%text for constrained optimization
\newcommand{\subjectto}{\textrm{subject to}\;}

%%% Local Variables:
%%% mode: latex
%%% End:

%memberships
\newcommand{\iV}[1][]{{i \in V_{#1}}}
\newcommand{\ijE}[1][]{(i,j) \in E_{#1}}

%operators
\DeclareMathOperator{\dg}{{deg}}
\DeclareMathOperator{\diam}{{Diam}}

%%% Local Variables: 
%%% mode: latex
%%% End: 

% This file was generated by the scriptgenerateNotation
% Do not modify this file directly

% Shortand notation for vectors and their derivatives
\providecommand{\va}{\vct{a}}
\providecommand{\dva}{\dot{\vct{a}}}
\providecommand{\tva}{\tilde{\vct{a}}}
\providecommand{\dtva}{\dot{\tilde{\vct{a}}}}
\providecommand{\vb}{\vct{b}}
\providecommand{\dvb}{\dot{\vct{b}}}
\providecommand{\tvb}{\tilde{\vct{b}}}
\providecommand{\dtvb}{\dot{\tilde{\vct{b}}}}
\providecommand{\vc}{\vct{c}}
\providecommand{\dvc}{\dot{\vct{c}}}
\providecommand{\tvc}{\tilde{\vct{c}}}
\providecommand{\dtvc}{\dot{\tilde{\vct{c}}}}
\providecommand{\vd}{\vct{d}}
\providecommand{\dvd}{\dot{\vct{d}}}
\providecommand{\tvd}{\tilde{\vct{d}}}
\providecommand{\dtvd}{\dot{\tilde{\vct{d}}}}
\providecommand{\ve}{\vct{e}}
\providecommand{\dve}{\dot{\vct{e}}}
\providecommand{\tve}{\tilde{\vct{e}}}
\providecommand{\dtve}{\dot{\tilde{\vct{e}}}}
\providecommand{\vf}{\vct{f}}
\providecommand{\dvf}{\dot{\vct{f}}}
\providecommand{\tvf}{\tilde{\vct{f}}}
\providecommand{\dtvf}{\dot{\tilde{\vct{f}}}}
\providecommand{\vg}{\vct{g}}
\providecommand{\dvg}{\dot{\vct{g}}}
\providecommand{\tvg}{\tilde{\vct{g}}}
\providecommand{\dtvg}{\dot{\tilde{\vct{g}}}}
\providecommand{\vh}{\vct{h}}
\providecommand{\dvh}{\dot{\vct{h}}}
\providecommand{\tvh}{\tilde{\vct{h}}}
\providecommand{\dtvh}{\dot{\tilde{\vct{h}}}}
\providecommand{\vi}{\vct{i}}
\providecommand{\dvi}{\dot{\vct{i}}}
\providecommand{\tvi}{\tilde{\vct{i}}}
\providecommand{\dtvi}{\dot{\tilde{\vct{i}}}}
\providecommand{\vj}{\vct{j}}
\providecommand{\dvj}{\dot{\vct{j}}}
\providecommand{\tvj}{\tilde{\vct{j}}}
\providecommand{\dtvj}{\dot{\tilde{\vct{j}}}}
\providecommand{\vk}{\vct{k}}
\providecommand{\dvk}{\dot{\vct{k}}}
\providecommand{\tvk}{\tilde{\vct{k}}}
\providecommand{\dtvk}{\dot{\tilde{\vct{k}}}}
\providecommand{\vl}{\vct{l}}
\providecommand{\dvl}{\dot{\vct{l}}}
\providecommand{\tvl}{\tilde{\vct{l}}}
\providecommand{\dtvl}{\dot{\tilde{\vct{l}}}}
\providecommand{\vm}{\vct{m}}
\providecommand{\dvm}{\dot{\vct{m}}}
\providecommand{\tvm}{\tilde{\vct{m}}}
\providecommand{\dtvm}{\dot{\tilde{\vct{m}}}}
\providecommand{\vn}{\vct{n}}
\providecommand{\dvn}{\dot{\vct{n}}}
\providecommand{\tvn}{\tilde{\vct{n}}}
\providecommand{\dtvn}{\dot{\tilde{\vct{n}}}}
\providecommand{\vo}{\vct{o}}
\providecommand{\dvo}{\dot{\vct{o}}}
\providecommand{\tvo}{\tilde{\vct{o}}}
\providecommand{\dtvo}{\dot{\tilde{\vct{o}}}}
\providecommand{\vp}{\vct{p}}
\providecommand{\dvp}{\dot{\vct{p}}}
\providecommand{\tvp}{\tilde{\vct{p}}}
\providecommand{\dtvp}{\dot{\tilde{\vct{p}}}}
\providecommand{\vq}{\vct{q}}
\providecommand{\dvq}{\dot{\vct{q}}}
\providecommand{\tvq}{\tilde{\vct{q}}}
\providecommand{\dtvq}{\dot{\tilde{\vct{q}}}}
\providecommand{\vr}{\vct{r}}
\providecommand{\dvr}{\dot{\vct{r}}}
\providecommand{\tvr}{\tilde{\vct{r}}}
\providecommand{\dtvr}{\dot{\tilde{\vct{r}}}}
\providecommand{\vs}{\vct{s}}
\providecommand{\dvs}{\dot{\vct{s}}}
\providecommand{\tvs}{\tilde{\vct{s}}}
\providecommand{\dtvs}{\dot{\tilde{\vct{s}}}}
\providecommand{\vt}{\vct{t}}
\providecommand{\dvt}{\dot{\vct{t}}}
\providecommand{\tvt}{\tilde{\vct{t}}}
\providecommand{\dtvt}{\dot{\tilde{\vct{t}}}}
\providecommand{\vu}{\vct{u}}
\providecommand{\dvu}{\dot{\vct{u}}}
\providecommand{\tvu}{\tilde{\vct{u}}}
\providecommand{\dtvu}{\dot{\tilde{\vct{u}}}}
\providecommand{\vv}{\vct{v}}
\providecommand{\dvv}{\dot{\vct{v}}}
\providecommand{\tvv}{\tilde{\vct{v}}}
\providecommand{\dtvv}{\dot{\tilde{\vct{v}}}}
\providecommand{\vw}{\vct{w}}
\providecommand{\dvw}{\dot{\vct{w}}}
\providecommand{\tvw}{\tilde{\vct{w}}}
\providecommand{\dtvw}{\dot{\tilde{\vct{w}}}}
\providecommand{\vx}{\vct{x}}
\providecommand{\dvx}{\dot{\vct{x}}}
\providecommand{\tvx}{\tilde{\vct{x}}}
\providecommand{\dtvx}{\dot{\tilde{\vct{x}}}}
\providecommand{\vy}{\vct{y}}
\providecommand{\dvy}{\dot{\vct{y}}}
\providecommand{\tvy}{\tilde{\vct{y}}}
\providecommand{\dtvy}{\dot{\tilde{\vct{y}}}}
\providecommand{\vz}{\vct{z}}
\providecommand{\dvz}{\dot{\vct{z}}}
\providecommand{\tvz}{\tilde{\vct{z}}}
\providecommand{\dtvz}{\dot{\tilde{\vct{z}}}}
\providecommand{\vA}{\vct{A}}
\providecommand{\dvA}{\dot{\vct{A}}}
\providecommand{\tvA}{\tilde{\vct{A}}}
\providecommand{\dtvA}{\dot{\tilde{\vct{A}}}}
\providecommand{\vB}{\vct{B}}
\providecommand{\dvB}{\dot{\vct{B}}}
\providecommand{\tvB}{\tilde{\vct{B}}}
\providecommand{\dtvB}{\dot{\tilde{\vct{B}}}}
\providecommand{\vC}{\vct{C}}
\providecommand{\dvC}{\dot{\vct{C}}}
\providecommand{\tvC}{\tilde{\vct{C}}}
\providecommand{\dtvC}{\dot{\tilde{\vct{C}}}}
\providecommand{\vD}{\vct{D}}
\providecommand{\dvD}{\dot{\vct{D}}}
\providecommand{\tvD}{\tilde{\vct{D}}}
\providecommand{\dtvD}{\dot{\tilde{\vct{D}}}}
\providecommand{\vE}{\vct{E}}
\providecommand{\dvE}{\dot{\vct{E}}}
\providecommand{\tvE}{\tilde{\vct{E}}}
\providecommand{\dtvE}{\dot{\tilde{\vct{E}}}}
\providecommand{\vF}{\vct{F}}
\providecommand{\dvF}{\dot{\vct{F}}}
\providecommand{\tvF}{\tilde{\vct{F}}}
\providecommand{\dtvF}{\dot{\tilde{\vct{F}}}}
\providecommand{\vG}{\vct{G}}
\providecommand{\dvG}{\dot{\vct{G}}}
\providecommand{\tvG}{\tilde{\vct{G}}}
\providecommand{\dtvG}{\dot{\tilde{\vct{G}}}}
\providecommand{\vH}{\vct{H}}
\providecommand{\dvH}{\dot{\vct{H}}}
\providecommand{\tvH}{\tilde{\vct{H}}}
\providecommand{\dtvH}{\dot{\tilde{\vct{H}}}}
\providecommand{\vI}{\vct{I}}
\providecommand{\dvI}{\dot{\vct{I}}}
\providecommand{\tvI}{\tilde{\vct{I}}}
\providecommand{\dtvI}{\dot{\tilde{\vct{I}}}}
\providecommand{\vJ}{\vct{J}}
\providecommand{\dvJ}{\dot{\vct{J}}}
\providecommand{\tvJ}{\tilde{\vct{J}}}
\providecommand{\dtvJ}{\dot{\tilde{\vct{J}}}}
\providecommand{\vK}{\vct{K}}
\providecommand{\dvK}{\dot{\vct{K}}}
\providecommand{\tvK}{\tilde{\vct{K}}}
\providecommand{\dtvK}{\dot{\tilde{\vct{K}}}}
\providecommand{\vL}{\vct{L}}
\providecommand{\dvL}{\dot{\vct{L}}}
\providecommand{\tvL}{\tilde{\vct{L}}}
\providecommand{\dtvL}{\dot{\tilde{\vct{L}}}}
\providecommand{\vM}{\vct{M}}
\providecommand{\dvM}{\dot{\vct{M}}}
\providecommand{\tvM}{\tilde{\vct{M}}}
\providecommand{\dtvM}{\dot{\tilde{\vct{M}}}}
\providecommand{\vN}{\vct{N}}
\providecommand{\dvN}{\dot{\vct{N}}}
\providecommand{\tvN}{\tilde{\vct{N}}}
\providecommand{\dtvN}{\dot{\tilde{\vct{N}}}}
\providecommand{\vO}{\vct{O}}
\providecommand{\dvO}{\dot{\vct{O}}}
\providecommand{\tvO}{\tilde{\vct{O}}}
\providecommand{\dtvO}{\dot{\tilde{\vct{O}}}}
\providecommand{\vP}{\vct{P}}
\providecommand{\dvP}{\dot{\vct{P}}}
\providecommand{\tvP}{\tilde{\vct{P}}}
\providecommand{\dtvP}{\dot{\tilde{\vct{P}}}}
\providecommand{\vQ}{\vct{Q}}
\providecommand{\dvQ}{\dot{\vct{Q}}}
\providecommand{\tvQ}{\tilde{\vct{Q}}}
\providecommand{\dtvQ}{\dot{\tilde{\vct{Q}}}}
\providecommand{\vR}{\vct{R}}
\providecommand{\dvR}{\dot{\vct{R}}}
\providecommand{\tvR}{\tilde{\vct{R}}}
\providecommand{\dtvR}{\dot{\tilde{\vct{R}}}}
\providecommand{\vS}{\vct{S}}
\providecommand{\dvS}{\dot{\vct{S}}}
\providecommand{\tvS}{\tilde{\vct{S}}}
\providecommand{\dtvS}{\dot{\tilde{\vct{S}}}}
\providecommand{\vT}{\vct{T}}
\providecommand{\dvT}{\dot{\vct{T}}}
\providecommand{\tvT}{\tilde{\vct{T}}}
\providecommand{\dtvT}{\dot{\tilde{\vct{T}}}}
\providecommand{\vU}{\vct{U}}
\providecommand{\dvU}{\dot{\vct{U}}}
\providecommand{\tvU}{\tilde{\vct{U}}}
\providecommand{\dtvU}{\dot{\tilde{\vct{U}}}}
\providecommand{\vV}{\vct{V}}
\providecommand{\dvV}{\dot{\vct{V}}}
\providecommand{\tvV}{\tilde{\vct{V}}}
\providecommand{\dtvV}{\dot{\tilde{\vct{V}}}}
\providecommand{\vW}{\vct{W}}
\providecommand{\dvW}{\dot{\vct{W}}}
\providecommand{\tvW}{\tilde{\vct{W}}}
\providecommand{\dtvW}{\dot{\tilde{\vct{W}}}}
\providecommand{\vX}{\vct{X}}
\providecommand{\dvX}{\dot{\vct{X}}}
\providecommand{\tvX}{\tilde{\vct{X}}}
\providecommand{\dtvX}{\dot{\tilde{\vct{X}}}}
\providecommand{\vY}{\vct{Y}}
\providecommand{\dvY}{\dot{\vct{Y}}}
\providecommand{\tvY}{\tilde{\vct{Y}}}
\providecommand{\dtvY}{\dot{\tilde{\vct{Y}}}}
\providecommand{\vZ}{\vct{Z}}
\providecommand{\dvZ}{\dot{\vct{Z}}}
\providecommand{\tvZ}{\tilde{\vct{Z}}}
\providecommand{\dtvZ}{\dot{\tilde{\vct{Z}}}}

\providecommand{\vbeta}{\boldsymbol{\beta}}
\providecommand{\dvbeta}{\dot{\boldsymbol{\beta}}}
\providecommand{\vsigma}{\boldsymbol{\sigma}}
\providecommand{\dvsigma}{\dot{\boldsymbol{\sigma}}}

\providecommand{\te}{\tilde{e}}

% Shortand notation for derivatives and bold of symbols
\providecommand{\dgamma}{\dot{\gamma}}
\providecommand{\ddgamma}{\ddot{\gamma}}
\providecommand{\vgamma}{\bm{\gamma}}
\providecommand{\dDelta}{\dot{\Delta}}
\providecommand{\ddDelta}{\ddot{\Delta}}
\providecommand{\vDelta}{\bm{\Delta}}
\providecommand{\dpi}{\dot{\pi}}
\providecommand{\ddpi}{\ddot{\pi}}
\providecommand{\vpi}{\bm{\pi}}
\providecommand{\dvarphi}{\dot{\varphi}}
\providecommand{\ddvarphi}{\ddot{\varphi}}
\providecommand{\vvarphi}{\bm{\varphi}}

% Shortand notation for matrices
\providecommand{\mA}{\vct{A}}
\providecommand{\mB}{\vct{B}}
\providecommand{\mC}{\vct{C}}
\providecommand{\mD}{\vct{D}}
\providecommand{\mE}{\vct{E}}
\providecommand{\mF}{\vct{F}}
\providecommand{\mG}{\vct{G}}
\providecommand{\mH}{\vct{H}}
\providecommand{\mI}{\vct{I}}
\providecommand{\mJ}{\vct{J}}
\providecommand{\mK}{\vct{K}}
\providecommand{\mL}{\vct{L}}
\providecommand{\mM}{\vct{M}}
\providecommand{\mN}{\vct{N}}
\providecommand{\mO}{\vct{O}}
\providecommand{\mP}{\vct{P}}
\providecommand{\mQ}{\vct{Q}}
\providecommand{\mR}{\vct{R}}
\providecommand{\mS}{\vct{S}}
\providecommand{\mT}{\vct{T}}
\providecommand{\mU}{\vct{U}}
\providecommand{\mV}{\vct{V}}
\providecommand{\mW}{\vct{W}}
\providecommand{\mX}{\vct{X}}
\providecommand{\mY}{\vct{Y}}
\providecommand{\mZ}{\vct{Z}}

% Shortand notation for calligraphic upper case letters
\providecommand{\cA}{\mathcal{A}}
\providecommand{\cB}{\mathcal{B}}
\providecommand{\cC}{\mathcal{C}}
\providecommand{\cD}{\mathcal{D}}
\providecommand{\cE}{\mathcal{E}}
\providecommand{\cF}{\mathcal{F}}
\providecommand{\cG}{\mathcal{G}}
\providecommand{\cH}{\mathcal{H}}
\providecommand{\cI}{\mathcal{I}}
\providecommand{\cJ}{\mathcal{J}}
\providecommand{\cK}{\mathcal{K}}
\providecommand{\cL}{\mathcal{L}}
\providecommand{\cM}{\mathcal{M}}
\providecommand{\cN}{\mathcal{N}}
\providecommand{\cO}{\mathcal{O}}
\providecommand{\cP}{\mathcal{P}}
\providecommand{\cQ}{\mathcal{Q}}
\providecommand{\cR}{\mathcal{R}}
\providecommand{\cS}{\mathcal{S}}
\providecommand{\cT}{\mathcal{T}}
\providecommand{\cU}{\mathcal{U}}
\providecommand{\cV}{\mathcal{V}}
\providecommand{\cW}{\mathcal{W}}
\providecommand{\cX}{\mathcal{X}}
\providecommand{\cY}{\mathcal{Y}}
\providecommand{\cZ}{\mathcal{Z}}

% Shortand notation for some tilded symbols and their derivatives
\providecommand{\tx}{\tilde{x}}
\providecommand{\dtx}{\dot{\tilde{x}}}
\providecommand{\ty}{\tilde{y}}
\providecommand{\dty}{\dot{\tilde{y}}}
\providecommand{\tvarphi}{\tilde{\varphi}}
\providecommand{\dtvarphi}{\dot{\tilde{\varphi}}}


%command for units of measure
\usepackage{units}

%S.I. units for some standard quantities
\newcommand{\upos}[1][]{\unit[#1]{m}}
\newcommand{\urad}[1][]{\unit[#1]{rad}}
\newcommand{\uvel}[1][]{\unitfrac[#1]{m}{s}}
\newcommand{\uacc}[1][]{\unitfrac[#1]{m}{s^2}}
\newcommand{\uavel}[1][]{\unitfrac[#1]{rad}{s}}
\newcommand{\uaacc}[1][]{\unitfrac[#1]{rad}{s^2}}
\newcommand{\uforce}[1][]{\unit[#1]{N}}
\newcommand{\utime}[1][]{\unit[#1]{s}}
\newcommand{\umass}[1][]{\unit[#1]{kg}}
\newcommand{\uspring}[1][]{\unitfrac[#1]{N}{m}}
\newcommand{\udamping}[1][]{\unitfrac[#1]{N\cdot s}{m}}
\newcommand{\ulength}[1][]{\unit[#1]{m}}
\newcommand{\uinertia}[1][]{\unit[#1]{kg\cdot m^2}}

%

%macro to define other macros for block-colored labels
% \newcommand{\newcolorlabel}[2]{%
%   \expandafter\newcommand\csname #1\endcsname[1]{%
%     \colorbox{#2}{\color{white}\textsf{\textbf{##1}}}}%
% }
  \newcommand{\newcolorlabel}[2]{%
  \expandafter\newcommand\csname #1\endcsname[1]{%
    \tikz[baseline]{\node[text=white,fill=#2,anchor=base,text height=1.3ex,text depth=0.1ex,font=\sffamily\bfseries]{##1}}}%
}

%macro to define other macros for comments 
%
\newcommand{\newcommenter}[2]{%
  \expandafter\newcommand\csname #1\endcsname[1]{%
    \fcolorbox{#2}{#2}{\color{white}\textsf{\textbf{#1}}}
    {\color{#2}##1}}%
  % comment to mention commenter
  \expandafter\newcommand\csname at#1\endcsname{%
    \fcolorbox{#2}{#2}{\color{white}\textsf{\textbf{@#1}}}
    {\color{#2}}}%
  % citation placeholder
  \expandafter\newcommand\csname #1cite\endcsname[1]{%
    \csname #1\endcsname {[##1]}
  }%
  % internal reference placeholder
  \expandafter\newcommand\csname #1ref\endcsname[1]{%
    \csname #1\endcsname {$\blacktriangleright$##1}
  }%
  % comment to highlight
  \expandafter\newcommand\csname #1hl\endcsname[2]{%
    \colorbox{#2}{\color{white}\textsf{\textbf{#1}}}\sethlcolor{Azure2}\hl{##2}~%
    \expandafter\ifx\csname commentarrow\endcsname\relax$\leftarrow$\else \commentarrow[#2]\fi~%
    {\color{#2}##1}}%
  % comment to strikeout
  \expandafter\newcommand\csname #1st\endcsname[2]{%
    \colorbox{#2}{\color{white}\textsf{\textbf{#1}}}\sout{##2}~%
    \expandafter\ifx\csname commentarrow\endcsname\relax$\leftarrow$\else \commentarrow[#2]\fi~%
    {\color{#2}##1}}%
}
% examples of the macro above
\newcommenter{TODO}{DodgerBlue1}
\newcommenter{rtron}{Green3}

%side review pointer
\newcommand{\toreview}{\tikz[overlay, remember picture]{\coordinate (center); \node[fill=red] at (current page.center |- center){};}}

%introduce the comment environment
\usepackage{comment}

%enable pdf annotation
\usepackage{pdfcomment}

%enable highlights
\usepackage{soul}

%enable strikeout text with the command \sout{}
\usepackage[normalem]{ulem}

%package for displayed text
\usepackage{csquotes}

%markup
\newcommand{\file}[1]{\colorbox{Ivory2!70!DarkSeaGreen2}{\texttt{#1}}}
\newcommand{\var}[1]{\colorbox{Ivory2}{\texttt{#1}}}

\newcommand{\function}[3][]{\ifnotempty{#1}{[\var{#1}]=}\file{\mbox{#2}}(\var{#3})}

\newcommand{\key}[1]{\mbox{\textcolor{DodgerBlue3}{\texttt{#1}}}}

\newcommand{\displayfunction}[3][]{\begin{displayquote}\function[#1]{#2}{#3}\end{displayquote}}
\newcommand{\vardim}[2]{[\texttt{#1}~$\times$~\texttt{#2}]}

% Allow easy definition of starred version of commands
% Ref: https://tex.stackexchange.com/questions/202504/macro-to-add-starred-version-of-command
\usepackage{suffix}

% Allow definition of environments with extra final code
\usepackage{environ}

% Define warning sign by extracting it from the Fourier package
% Source: https://tex.stackexchange.com/questions/159669/how-to-print-a-warning-sign-triangle-with-exclamation-point
\newcommand{\danger}{\textcolor{red}{\fontencoding{U}\fontfamily{futs}\selectfont\char49\relax}}

% Insert a prefix-argument-postfix text only if argument is non-empty
% Needs to use a savebox to avoid evaluating the argument multiple times
\makeatletter
\newsavebox{\boxifnotempty}
\newcommand{\displayifnotempty}[3]{\sbox\boxifnotempty{#2}\setbox0=\hbox{\usebox{\boxifnotempty}\unskip}%
  \ifdim\wd0=0pt
  \else
  #1\usebox{\boxifnotempty}#3%
  \fi%
}

% \ifempty{1}{2} displays 2 if 1 is empty
\newcommand{\ifempty}[2]{\setbox0=\hbox{#1\unskip}%
  \ifdim\wd0=0pt%
  #2%
  \fi%
}

% \ifnotempty{1}{2} displays 2 if 1 is not empty
\newcommand{\ifnotempty}[2]{\setbox0=\hbox{#1\unskip}%
  \ifdim\wd0>0pt%
  #2%
  \fi%
}
\makeatother

% \switchempty{1}{2}{3} displays 2 if 1 is empty, otherwise 3
\newcommand{\switchifempty}[3]{\sbox\boxifnotempty{#1}\setbox0=\hbox{\usebox{\boxifnotempty}\unskip}%
  \ifdim\wd0=0pt{}%
  #2%
  \else{}%
  #3%
  \usebox{\boxifnotempty}%
  \fi%
}


% introduce the algorithmic environment and the algorithm floats
\makeatletter
\@ifundefined{chapter}{\usepackage{algorithm}}{\usepackage[chapter]{algorithm}}
\makeatother
\usepackage{algorithmicx}
\usepackage{algpseudocode}
\makeatother%

% macros for storing definitions across compilations
\usepackage{scrlfile}

\makeatletter
%mark a definition to be stored in the aux file
\newcommand*\newstoreddef[1]{
  \BeforeClosingMainAux{%
    \immediate\write\@auxout{%
      \string\restoredef{#1}{\csname #1\endcsname}%
    }%
  }%
}
%used by the aux file to restore the definition
\newcommand*{\restoredef}[2]{% used at the aux file
  \expandafter\gdef\csname stored@#1\endcsname{#2}%
}
%show the stored definition (user command to ask for the value)
\newcommand*{\storeddef}[1]{
  \@ifundefined{stored@#1}{0}{\csname stored@#1\endcsname}%
}
\makeatother

%Add values to non-counter definitions (works with non-integers)
\newcommand{\addtovar}[2]{\pgfmathparse{#1+#2}\xdef#1{\pgfmathresult}}
\newcommand{\zerovar}[1]{\xdef#1{0}}

%Insert content of a PGF variable
\newcommand{\pgfprint}[1]{\pgfmathparse{#1}\pgfmathresult}

%Package to get PDF page numbers
\usepackage{pageslts}
\pagenumbering{arabic}

%Output content of enviroment both to the document and to the log file
%In the log file, the content is marked by start/end delimiters, and
%macros are not expanded.
\NewEnviron{tee}{\BODY\typeout{Marker Tee [start] ^^J \BODY ^^JMaker Tee [end]}}

% Clever cross-references
\usepackage{cleveref}


%%% Local Variables:
%%% mode: latex
%%% End:


%%% Local Variables:
%%% mode: latex
%%% End:

%TikZ and common libraries
\usepackage{tikz}
\usetikzlibrary{calc}
\usetikzlibrary{matrix}
\usetikzlibrary{chains,scopes}
\usetikzlibrary{shapes.geometric}
\usetikzlibrary{arrows.meta}
\usetikzlibrary{decorations.markings}
\usetikzlibrary{decorations.pathreplacing}
\usetikzlibrary{backgrounds}

%Example figure
\newcommand{\picplaceholder}[1]{\tikz{\draw (-2,-1.5) rectangle (2,1.5); \node[font=\tiny] {#1};}}

%Draw normalized vector between two coordinates
\newcommand\normalize[2][(0,0)]{%
  \draw[blue,->] (#1) -- ($(#1)!1cm!(#2)$) coordinate (#1#2norm)}

%Quotatures
\tikzset{
  dim above/.style={to path={\pgfextra{
        \pgfinterruptpath
        \draw[>=latex,|->|] let
        \p1=($(\tikztostart)!1.5em!90:(\tikztotarget)$),
        \p2=($(\tikztotarget)!1.5em!-90:(\tikztostart)$)
        in(\p1) -- (\p2) node[pos=.5,sloped,above]{#1};
        \endpgfinterruptpath
      }
    }
  },
  dim double above/.style={to path={\pgfextra{
        \pgfinterruptpath
        \draw[>=latex,|->|] let
        \p1=($(\tikztostart)!3em!90:(\tikztotarget)$),
        \p2=($(\tikztotarget)!3em!-90:(\tikztostart)$)
        in(\p1) -- (\p2) node[pos=.5,sloped,above]{#1};
        \endpgfinterruptpath
      }
    }
  },
  dim below/.style={to path={\pgfextra{
        \pgfinterruptpath
        \draw[>=latex,|->|] let 
        \p1=($(\tikztostart)!-1em!-90:(\tikztotarget)$),
        \p2=($(\tikztotarget)!-1em!90:(\tikztostart)$)
        in (\p1) -- (\p2) node[pos=.5,sloped,below]{#1};
        \endpgfinterruptpath
      }
    }
  },
}

%Right angle symbol
\tikzset{
    right angle quadrant/.code={
        \pgfmathsetmacro\quadranta{{1,1,-1,-1}[#1-1]}     % Arrays for selecting quadrant
        \pgfmathsetmacro\quadrantb{{1,-1,-1,1}[#1-1]}},
    right angle quadrant=1, % Make sure it is set, even if not called explicitly
    right angle length/.code={\def\rightanglelength{#1}},   % Length of symbol
    right angle length=2ex, % Make sure it is set...
    right angle symbol/.style n args={3}{
        insert path={
            let \p0 = ($(#1)!(#3)!(#2)$) in     % Intersection
                let \p1 = ($(\p0)!\quadranta*\rightanglelength!(#3)$), % Point on base line
                \p2 = ($(\p0)!\quadrantb*\rightanglelength!(#2)$) in % Point on perpendicular line
                let \p3 = ($(\p1)+(\p2)-(\p0)$) in  % Corner point of symbol
            (\p1) -- (\p3) -- (\p2)
        }
    }
}

%Horizontally fit an image between two coordinates
\newcommand{\imageBetween}[3]{\draw (#2) let \p1 = ($ (#3) - (#2) $) in node[inner sep=0pt,anchor=south west, text width={veclen(\x1,\y1)}] {\includegraphics[width=\linewidth]{#1}};}

%Get angle between a line going through two points and the horizontal
%direction
\newcommand{\pgfextractangle}[3]{%
    \pgfmathanglebetweenpoints{\pgfpointanchor{#2}{center}}
                              {\pgfpointanchor{#3}{center}}
    \global\let#1\pgfmathresult  
}

%Arrow to be used to indicate something in the text
\usetikzlibrary{shapes.arrows}
\newcommand{\commentarrow}[1][Azure4]{\tikz[baseline=-3pt]{\node[shape border uses incircle, fill=#1,rotate=180,single arrow, inner sep=1pt, minimum size=6pt, single arrow head extend=2pt]{};}}

%Draw a test grid to help with positioning
\newcommand{\tikztestgrid}[3][cm]{
  \coordinate (testgrid origin) at (0,0);
  \coordinate (testgrid bl) at (-#2#1/2,-#3#1/2);
  \coordinate (testgrid ur) at (#2#1/2,#3#1/2);
  \begin{scope}[opacity=0.5]
    \draw[cyan!50] (testgrid bl) grid[step=0.1#1] (testgrid ur);
    \draw[cyan] (testgrid bl) grid[step=1#1] (testgrid ur);
  \end{scope}
  \draw[red,-latex] (testgrid bl|-testgrid origin) -- (testgrid ur|-testgrid origin);
  \draw[red,-latex] (testgrid bl-|testgrid origin) -- (testgrid ur-|testgrid origin);
  \node[circle,fill=black,minimum size=3pt,inner sep=0pt] at (testgrid origin) {};
}

\tikzset{ax/.style={-latex,line width=2pt}}
\newcommand{\ax}[5][]{\begin{scope}[xshift=#3,yshift=#4,rotate=#5]%
    \node at (-0.3,-0.3) {$#1$};%
    \coordinate (ax#2o) at (0,0);%
    \draw[ax] (0,0) -- (1,0) node[pos=1.2] (ax#2a) {};%
    \draw[ax] (0,0) -- (0,1) node[pos=1.2] (ax#2b) {};%
  \end{scope}}

\newcommand{\axlabels}[3]{\node at (ax#1a) {$#2$};
  \node at (ax#1b) {$#3$};
}
\newcommand{\axup}[3][0]{\begin{scope}[rotate=#1]
\fill[ax] (ax#2o) circle[radius=3pt];
\path (ax#2o) ++(0.25,-0.25) node{$#3$};
\end{scope}
}
\newcommand{\axdown}[3]{\begin{scope}[rotate=#3]
    \draw[line width=2pt] (ax#1o) +(-3pt,-3pt) -- +(3pt,3pt) +(-3pt,3pt) -- +(3pt,-3pt);
    \path (ax#1o) ++(0.25,-0.25) node{$#2$};
  \end{scope}
}

\newcommand{\gpoint}[3][0]{\fill[ax] (#2) circle[radius=3pt];
  \path (#2) ++(#1:0.4) node[ax]{$#3$};
}

\tikzset{camera/.style={fill=Sienna1,fill opacity=0.5},%
image plane/.style={draw=RoyalBlue3,line width=2pt}}
\newcommand{\camera}[5][]{\begin{scope}[xshift=#3,yshift=#4,rotate=#5]
\ax[#1]{#2}{0}{0}{0}
\draw[image plane] (-1.2,1) -- (1.2,1);
\draw[camera] (0,0) -- (0.5,0.5) -- (-0.5,0.5) -- cycle; 
\end{scope}}

\newcommand{\gridworld}{\draw[gray] (0,0) grid[step=1cm] (5,5);}

\newcommand{\wframe}{\prescript{w}{}}
\newcommand{\Wframe}{\prescript{\cW}{}}
\newcommand{\iframe}{\prescript{i}{}}
\newcommand{\jframe}{\prescript{j}{}}
\newcommand{\bframe}[1][]{\prescript{b_{#1}}{}}
\newcommand{\Bframe}[1][]{\prescript{\cB_{#1}}{}}

\newcommand{\BtoWpose}{(\Wframe R_\cB,\Wframe T_\cB)}
\newcommand{\BtoWrot}{\Wframe R_\cB}
\newcommand{\BtoWtransl}{\Wframe T_\cB}
\newcommand{\BtoWposeVel}{(\Wframe{}\dot{R}_\cB,\Wframe{}\dot{T}_\cB)}
\newcommand{\BtoWrotVel}{\Wframe{}\dot{R}_\cB}
\newcommand{\BtoWtranslVel}{\Wframe{}\dot{T}_\cB}

\newcommand{\WtoBpose}{(\Bframe R_\cW,\Bframe T_\cW)}
\newcommand{\WtoBrot}{\Bframe R_\cW}
\newcommand{\WtoBtransl}{\Bframe T_\cW}
\newcommand{\WtoBposeVel}{(\Bframe{}\dot{R}_\cW,\Bframe{}\dot{T}_\cW)}
\newcommand{\WtoBrotVel}{\Bframe{}\dot{R}_\cW}
\newcommand{\WtoBtranslVel}{\Bframe{}\dot{T}_\cW}

\newcommand{\Bframex}{\Bframe x}
\newcommand{\Wframex}{\Wframe x}

% shorthand notation for 2-D rotation written as function of theta
\newcommand{\rottheta}{\left[ \begin{smallmatrix}
    \cos(\theta) & -\sin(\theta)\\ \sin(\theta) & \cos(\theta)
  \end{smallmatrix}
\right]}

%names of path planning algorithms
\newcommand{\rrt}{$\mathtt{RRT}$}
\newcommand{\rrtstar}{$\mathtt{RRT^*}$}
\newcommand{\astar}{$\mathtt{A^*}$}


%%% Local Variables: 
%%% mode: latex
%%% End: 

% Make font for captions smaller (saves space and visually separates them from the main text)
\usepackage[font={footnotesize,sf}]{caption}

% Print menu and key command descriptions
% \usepackage{menukeys}

% make underscore work without escaping in text mode
\usepackage{underscore}

% Common style that will be used for titles
\newcommand{\titlestyle}{\color{DodgerBlue3}\sffamily}

% Perfect 1-inch margins
\usepackage[margin=1in]{geometry}
\usepackage[medium,compact]{titlesec}
\usepackage{longtable}
\titleformat*{\section}{\titlestyle \Large}
\titleformat*{\subsection}{\titlestyle \large \itshape}

% Lists (requires enumitem package):
% Font for description
\setlist[description]{font=\color{DodgerBlue3}}
% Compact lists
\setlist[1]{\labelindent=\parindent}
\setlist[1]{topsep=3pt, itemsep=0pt}

% Format main title
\let\prevtitle\title
\renewcommand{\title}[1]{\prevtitle{\titlestyle #1}}
\let\prevauthor\author
\renewcommand{\author}[1]{\prevauthor{\titlestyle #1}}
\let\prevdate\date
\renewcommand{\date}[1]{\prevdate{\titlestyle #1}}

% Displayed quotations (used to define an environment for function specification)
\usepackage{csquotes}

% Include code
% Include code and org files
\usepackage{listings}
% Smaller fixed-width font, automatic breaking of long lines, ignore indentation in list environments
\lstset{basicstyle=\footnotesize\ttfamily,breaklines=true,resetmargins=true}


% formatting for the description of a function's input/output arguments
\mdfdefinestyle{functionspecification}{%
  linecolor=DodgerBlue3,
  rightline=false, topline=false, bottomline=false,
  outerlinewidth=1.35,
  roundcorner=10pt,
  innerleftmargin=1.2em,
  innertopmargin=0.3em,
  innerbottommargin=0.3em,
  leftmargin=1pt,
  skipabove=0.5em,
  skipbelow=0.5em,
  backgroundcolor=white}

\newcommand{\dstar}{$\bigstar$}

\usetikzlibrary{shapes.geometric}

%colors for the players
\definecolor{player1}{named}{RoyalBlue1}
\definecolor{player2}{named}{Firebrick1}
\definecolor{player3}{named}{Gold1}
\definecolor{player4}{named}{SpringGreen2}

%sizes of board elements
\newcommand{\placesize}{7mm}
\newcommand{\cardwidth}{3.33cm}
\newcommand{\cardheight}{\cardwidth*4/3}

%line and node styles for the board
\tikzset{link/.style={very thick},
  place/.style={circle,minimum size=\placesize,draw=black,line width=1mm},
  start/.style={fill=red,place},
  netnode/.style={fill=blue,place},
  polygon/.style={draw=orange,line width=2mm},
  card place/.style={polygon,align=center,minimum width=\cardwidth,minimum height=\cardheight},
}

%styles for pawn symbols
\tikzset{triang/.style={regular polygon, regular polygon sides=3},
  starplus/.style={star,star points=4,star point ratio=2.5}
}

\newcommand{\pawn}[3]{\node[circle,fill=player#1,minimum size=\placesize] at (#3) {};%
  \node[inner sep=0pt,fill=white,minimum size=0.8*\placesize,#2] at (#3) {};}



\newcommand{\gamename}{{\sffamily NetMatch: the Game}}

\title{\gamename{}  (version 0.1)}
\author{Roberto Tron}
\date{\today}
\begin{document}
\pagestyle{empty}

\maketitle

\section*{Overview}
The goal of \gamename{} is to discover your secret word by transferring and matching information across the network.
\section*{Game content}
\begin{itemize}
\item 64 feature pawns (4 copies of 4 symbols for each of the  4 colors).
\item 16 word tokens (4 symbols for each of the 4 colors).
\item 75 word cards.
\item 1 game board (with 1 dodecagonal network and 4 card boxes).
\end{itemize}
Note that all the pawns, tokens and cards need to be cut (using, e.g., scissors) from their respective sheets.

\section*{Rules}
\gamename{} is a game for 2-4 players. The rules below are for the 4 player version. See the corresponding section for the 3 players and 2 players variations.
\subsection*{Game preparation}
First, each player picks a color, and the deck of word cards is shuffled. Then, each player will setup the game for \emph{the next player on the right} using the following steps.
\begin{enumerate}
\item Pick a card and read the four-letter word on it, keeping it secret from other players.
\item Write each one of the letters under one of the \emph{other player's} letter token (not your own color).
\item Place the card (face down) in the \emph{other player's} card box on the board, and the four letter tokens on top of it.
\end{enumerate}
Finally, each player will place his own colored pawns on the sides of the dodecagonal network marked with their color; each side cannot contain the same symbol twice, and the symbols need to appear in the same order in each sector (see Figure~\ref{fig:pawn-arrangment}). Note that, in the 4 player game, only 48 of the pawns will be used.

\subsection*{Game play}
The youngest player goes first, and the turn is passed to the right.
\paragraph{Moves}
Each turn, the player can attempt the \emph{word gambit} (see below), or can move one of the pawns from a \emph{origin node} to a \emph{destination node} according to the following rules:
\begin{itemize}
\item The \emph{origin} and \emph{destination nodes} must be neighbors (i.e., directly connected with a black line).
\item The \emph{destination node} must be empty, or be occupied by a pawn with the same color and symbol.
\end{itemize}
If two pawns with the same color and symbol are in the same node, the player can declare a \emph{pairwise match}, and must remove one of the pawns involved from the board. If, by doing so, only one copy of a pawn remain in play (i.e., there have been two pairwise matches for that color/shape combination), the player can declare a \emph{multi-way match}, pick the corresponding letter token, and discover the letter written under it.
\paragraph{Winning}
A player can win by discovering all their four letters (by matching all the pawns and collecting all the tokens), which would allow them to read the word on their card.

Alternatively, a player can attempt, on his own turn, the \emph{word gambit} by stating a four letter word, and then turning their word card so that all players can see it: if the stated word and the word on the card match, then the player immediately wins; if they don't, the player immediately looses, and all their pawns are removed from the game.

\section*{Variations for two and three players}
There are two variations for playing the game with less than four players. In all cases, the game play is the same, only the game preparation is different.
\paragraph{Partial board variation} The same preparations as the four-players version are used, except that the colors that are not played are ignored. In this way, there are less pawns on the board, and there is less interaction between players.
\paragraph{Full board variation} For a three-player game, each player starts from four sides of the board instead of three. The sides should follow a cyclic alternating of colors (ignore the coloring already present on the board). For a two-player game, each player takes control of two colors; to win, a player must discover the words for both colors. With this variation, more matches are required to win.
\begin{figure}
\centering
\begin{tikzpicture}
  \begin{scope}[transform canvas={scale=0.2}]
  %This file has been generated by POCcostFunction.py

%Player sector colors
\path[fill=player1](8.8,0.0) -- (7.6210235533,4.4) -- (6.92820323028,4.0) -- (8.0,0.0) -- cycle;
\path[fill=player1](-4.4,7.6210235533) -- (-7.6210235533,4.4) -- (-6.92820323028,4.0) -- (-4.0,6.92820323028) -- cycle;
\path[fill=player1](-4.4,-7.6210235533) -- (-1.61653377487e-15,-8.8) -- (-1.46957615898e-15,-8.0) -- (-4.0,-6.92820323028) -- cycle;
\path[fill=player2](7.6210235533,4.4) -- (4.4,7.6210235533) -- (4.0,6.92820323028) -- (6.92820323028,4.0) -- cycle;
\path[fill=player2](-7.6210235533,4.4) -- (-8.8,1.07768918325e-15) -- (-8.0,9.79717439318e-16) -- (-6.92820323028,4.0) -- cycle;
\path[fill=player2](-1.61653377487e-15,-8.8) -- (4.4,-7.6210235533) -- (4.0,-6.92820323028) -- (-1.46957615898e-15,-8.0) -- cycle;
\path[fill=player3](4.4,7.6210235533) -- (5.38844591625e-16,8.8) -- (4.89858719659e-16,8.0) -- (4.0,6.92820323028) -- cycle;
\path[fill=player3](-8.8,1.07768918325e-15) -- (-7.6210235533,-4.4) -- (-6.92820323028,-4.0) -- (-8.0,9.79717439318e-16) -- cycle;
\path[fill=player3](4.4,-7.6210235533) -- (7.6210235533,-4.4) -- (6.92820323028,-4.0) -- (4.0,-6.92820323028) -- cycle;
\path[fill=player4](5.38844591625e-16,8.8) -- (-4.4,7.6210235533) -- (-4.0,6.92820323028) -- (4.89858719659e-16,8.0) -- cycle;
\path[fill=player4](-7.6210235533,-4.4) -- (-4.4,-7.6210235533) -- (-4.0,-6.92820323028) -- (-6.92820323028,-4.0) -- cycle;
\path[fill=player4](7.6210235533,-4.4) -- (8.8,0.0) -- (8.0,0.0) -- (6.92820323028,-4.0) -- cycle;
%Polygon
\draw[polygon] (8.0,0.0) -- (6.92820323028,4.0) -- (4.0,6.92820323028) -- (4.89858719659e-16,8.0) -- (-4.0,6.92820323028) -- (-6.92820323028,4.0) -- (-8.0,9.79717439318e-16) -- (-6.92820323028,-4.0) -- (-4.0,-6.92820323028) -- (-1.46957615898e-15,-8.0) -- (4.0,-6.92820323028) -- (6.92820323028,-4.0) -- cycle;
%Nodes and links
\path[start] (7.86602540378,0.5) coordinate[start] (s0) (7.59807621135,1.5) coordinate[start] (s1) (7.33012701892,2.5) coordinate[start] (s2) (7.06217782649,3.5) coordinate[start] (s3) (6.56217782649,4.36602540378) coordinate[start] (s4) (5.83012701892,5.09807621135) coordinate[start] (s5) (5.09807621135,5.83012701892) coordinate[start] (s6) (4.36602540378,6.56217782649) coordinate[start] (s7) (3.5,7.06217782649) coordinate[start] (s8) (2.5,7.33012701892) coordinate[start] (s9) (1.5,7.59807621135) coordinate[start] (s10) (0.5,7.86602540378) coordinate[start] (s11) (-0.5,7.86602540378) coordinate[start] (s12) (-1.5,7.59807621135) coordinate[start] (s13) (-2.5,7.33012701892) coordinate[start] (s14) (-3.5,7.06217782649) coordinate[start] (s15) (-4.36602540378,6.56217782649) coordinate[start] (s16) (-5.09807621135,5.83012701892) coordinate[start] (s17) (-5.83012701892,5.09807621135) coordinate[start] (s18) (-6.56217782649,4.36602540378) coordinate[start] (s19) (-7.06217782649,3.5) coordinate[start] (s20) (-7.33012701892,2.5) coordinate[start] (s21) (-7.59807621135,1.5) coordinate[start] (s22) (-7.86602540378,0.5) coordinate[start] (s23) (-7.86602540378,-0.5) coordinate[start] (s24) (-7.59807621135,-1.5) coordinate[start] (s25) (-7.33012701892,-2.5) coordinate[start] (s26) (-7.06217782649,-3.5) coordinate[start] (s27) (-6.56217782649,-4.36602540378) coordinate[start] (s28) (-5.83012701892,-5.09807621135) coordinate[start] (s29) (-5.09807621135,-5.83012701892) coordinate[start] (s30) (-4.36602540378,-6.56217782649) coordinate[start] (s31) (-3.5,-7.06217782649) coordinate[start] (s32) (-2.5,-7.33012701892) coordinate[start] (s33) (-1.5,-7.59807621135) coordinate[start] (s34) (-0.5,-7.86602540378) coordinate[start] (s35) (0.5,-7.86602540378) coordinate[start] (s36) (1.5,-7.59807621135) coordinate[start] (s37) (2.5,-7.33012701892) coordinate[start] (s38) (3.5,-7.06217782649) coordinate[start] (s39) (4.36602540378,-6.56217782649) coordinate[start] (s40) (5.09807621135,-5.83012701892) coordinate[start] (s41) (5.83012701892,-5.09807621135) coordinate[start] (s42) (6.56217782649,-4.36602540378) coordinate[start] (s43) (7.06217782649,-3.5) coordinate[start] (s44) (7.33012701892,-2.5) coordinate[start] (s45) (7.59807621135,-1.5) coordinate[start] (s46) (7.86602540378,-0.5) coordinate[start] (s47);
\path (-5.66458950116,-2.85371291234) coordinate[netnode] (n0) (-3.13953676463,-5.52085651601) coordinate[netnode] (n1) (2.28766112508,0.678831079224) coordinate[netnode] (n2) (5.70027506402,1.67571087853) coordinate[netnode] (n3) (0.162584255511,-2.39218233072) coordinate[netnode] (n4) (0.453712276129,-6.3681749908) coordinate[netnode] (n5) (-1.34171407598,5.95105537458) coordinate[netnode] (n6) (2.0758861129,-3.595689675) coordinate[netnode] (n7) (-6.04119876802,-1.06253756873) coordinate[netnode] (n8) (-1.37359793813,-5.95783698463) coordinate[netnode] (n9) (-2.05612586019,-3.73415858111) coordinate[netnode] (n10) (6.3509586276,-0.0416470987682) coordinate[netnode] (n11) (-4.30143454416,-4.09863362309) coordinate[netnode] (n12) (-2.25569745535,3.90707185545) coordinate[netnode] (n13) (-6.27206504806,0.823498351809) coordinate[netnode] (n14) (5.30372947143,3.47871124844) coordinate[netnode] (n15) (2.42292817113,5.84345970869) coordinate[netnode] (n16) (-3.81735339875,-1.91593227543) coordinate[netnode] (n17) (-1.73178248606,-1.64173770211) coordinate[netnode] (n18) (0.588091421565,6.10312705637) coordinate[netnode] (n19) (2.17912942633,-5.70194175276) coordinate[netnode] (n20) (2.03091755798,3.92679719021) coordinate[netnode] (n21) (3.8254204051,-2.28856060577) coordinate[netnode] (n22) (3.94086224942,4.70049587844) coordinate[netnode] (n23) (-5.57947267348,2.54232416471) coordinate[netnode] (n24) (0.00591003467172,-0.0102629739174) coordinate[netnode] (n25) (-0.124381953155,4.14846517207) coordinate[netnode] (n26) (5.28809431936,-3.57713120649) coordinate[netnode] (n27) (3.84836348643,-4.7382245695) coordinate[netnode] (n28) (-1.1949128307,2.06968815649) coordinate[netnode] (n29) (-3.53043878079,2.18202531352) coordinate[netnode] (n30) (0.916983837086,2.17886126415) coordinate[netnode] (n31) (3.56797064983,2.34784811551) coordinate[netnode] (n32) (4.26191591277,-0.0865436271976) coordinate[netnode] (n33) (-2.34545730164,0.295339841832) coordinate[netnode] (n34) (-3.20785167527,5.55624196324) coordinate[netnode] (n35) (5.84638346349,-1.78947744498) coordinate[netnode] (n36) (1.99038326562,-1.3369584776) coordinate[netnode] (n37) (-4.41618543933,0.204688715345) coordinate[netnode] (n38) (0.0690820574776,-4.4571725381) coordinate[netnode] (n39) (-4.48285212832,4.13754150742) coordinate[netnode] (n40);
\draw[link] (n38)  edge  (n14) (n38)  edge  (n17) (n38)  edge  (n24) (n38)  edge  (n34) (n39)  edge  (n4) (n39)  edge  (n5) (n39)  edge  (n7) (n39)  edge  (n10) (n39)  edge  (n18) (n30)  edge  (n13) (n30)  edge  (n24) (n30)  edge  (n29) (n30)  edge  (n34) (n30)  edge  (n38) (n30)  edge  (n40) (n31)  edge  (n2) (n31)  edge  (n21) (n31)  edge  (n25) (n31)  edge  (n26) (n31)  edge  (n29) (n31)  edge  (n32) (n32)  edge  (n3) (n32)  edge  (n15) (n32)  edge  (n33) (n33)  edge  (n3) (n33)  edge  (n11) (n33)  edge  (n37) (n34)  edge  (n17) (n34)  edge  (n25) (n35)  edge  (s14) (n35)  edge  (s15) (n7)  edge  (n37) (s44)  edge  (n27) (s44)  edge  (n36) (s45)  edge  (n27) (s45)  edge  (n36) (s46)  edge  (n11) (s46)  edge  (n36) (s47)  edge  (n11) (s47)  edge  (n36) (n12)  edge  (s28) (n12)  edge  (s29) (n12)  edge  (s30) (n12)  edge  (s31) (n12)  edge  (n10) (n12)  edge  (n17) (n13)  edge  (n6) (n13)  edge  (n26) (n13)  edge  (n29) (n13)  edge  (n35) (n13)  edge  (n40) (n10)  edge  (n4) (n10)  edge  (n17) (n10)  edge  (n18) (n11)  edge  (n36) (s7)  edge  (n16) (s24)  edge  (n14) (s4)  edge  (n15) (n18)  edge  (n4) (n18)  edge  (n17) (n18)  edge  (n25) (n18)  edge  (n34) (n19)  edge  (n26) (s19)  edge  (n24) (s18)  edge  (n35) (s13)  edge  (n6) (s13)  edge  (n19) (s12)  edge  (n6) (s12)  edge  (n19) (s11)  edge  (n16) (s11)  edge  (n19) (s10)  edge  (n16) (s10)  edge  (n19) (s17)  edge  (n35) (s16)  edge  (n35) (s38)  edge  (n5) (s35)  edge  (n5) (s35)  edge  (n9) (s34)  edge  (n5) (s34)  edge  (n9) (s37)  edge  (n5) (s37)  edge  (n20) (s36)  edge  (n5) (s36)  edge  (n20) (s22)  edge  (n14) (s22)  edge  (n24) (s9)  edge  (n16) (s9)  edge  (n19) (s8)  edge  (n16) (s23)  edge  (n14) (s3)  edge  (n3) (s3)  edge  (n15) (s2)  edge  (n3) (s2)  edge  (n15) (s1)  edge  (n3) (s1)  edge  (n11) (s0)  edge  (n3) (s0)  edge  (n11) (s6)  edge  (n15) (s5)  edge  (n15) (n40)  edge  (s16) (n40)  edge  (s17) (n40)  edge  (s18) (n40)  edge  (s19) (n29)  edge  (n25) (n29)  edge  (n26) (n29)  edge  (n34) (n28)  edge  (s39) (n28)  edge  (s40) (n28)  edge  (s41) (n28)  edge  (s42) (n28)  edge  (s43) (n28)  edge  (n7) (n23)  edge  (s4) (n23)  edge  (s5) (n23)  edge  (s6) (n23)  edge  (s7) (n23)  edge  (s8) (n23)  edge  (n21) (n22)  edge  (n7) (n22)  edge  (n27) (n22)  edge  (n33) (n22)  edge  (n36) (n22)  edge  (n37) (n21)  edge  (n16) (n21)  edge  (n26) (n21)  edge  (n32) (n20)  edge  (s38) (n20)  edge  (s39) (n20)  edge  (s40) (n20)  edge  (n7) (n27)  edge  (s41) (n27)  edge  (s42) (n27)  edge  (s43) (n26)  edge  (n6) (n25)  edge  (n37) (n8)  edge  (s23) (n8)  edge  (s24) (n8)  edge  (s25) (n8)  edge  (s26) (n8)  edge  (s27) (n8)  edge  (n38) (n9)  edge  (s32) (n9)  edge  (s33) (n9)  edge  (n39) (s20)  edge  (n14) (s20)  edge  (n24) (s21)  edge  (n14) (s21)  edge  (n24) (n0)  edge  (s25) (n0)  edge  (s26) (n0)  edge  (s27) (n0)  edge  (s28) (n0)  edge  (s29) (n0)  edge  (n17) (n1)  edge  (s30) (n1)  edge  (s31) (n1)  edge  (s32) (n1)  edge  (s33) (n1)  edge  (n9) (n1)  edge  (n10) (n2)  edge  (n21) (n2)  edge  (n25) (n2)  edge  (n32) (n2)  edge  (n33) (n2)  edge  (n37) (n4)  edge  (n7) (n4)  edge  (n25) (n4)  edge  (n37) (n6)  edge  (s14) (s15)  edge  (n6);
%Named coordinates for starting locations
\path coordinate (Spl1se0pa0) at (7.86602540378,0.5) coordinate (Spl1se0pa1) at (7.59807621135,1.5) coordinate (Spl1se0pa2) at (7.33012701892,2.5) coordinate (Spl1se0pa3) at (7.06217782649,3.5) ;
\path coordinate (Spl2se0pa0) at (6.56217782649,4.36602540378) coordinate (Spl2se0pa1) at (5.83012701892,5.09807621135) coordinate (Spl2se0pa2) at (5.09807621135,5.83012701892) coordinate (Spl2se0pa3) at (4.36602540378,6.56217782649) ;
\path coordinate (Spl3se0pa0) at (3.5,7.06217782649) coordinate (Spl3se0pa1) at (2.5,7.33012701892) coordinate (Spl3se0pa2) at (1.5,7.59807621135) coordinate (Spl3se0pa3) at (0.5,7.86602540378) ;
\path coordinate (Spl4se0pa0) at (-0.5,7.86602540378) coordinate (Spl4se0pa1) at (-1.5,7.59807621135) coordinate (Spl4se0pa2) at (-2.5,7.33012701892) coordinate (Spl4se0pa3) at (-3.5,7.06217782649) ;
\path coordinate (Spl1se1pa0) at (-4.36602540378,6.56217782649) coordinate (Spl1se1pa1) at (-5.09807621135,5.83012701892) coordinate (Spl1se1pa2) at (-5.83012701892,5.09807621135) coordinate (Spl1se1pa3) at (-6.56217782649,4.36602540378) ;
\path coordinate (Spl2se1pa0) at (-7.06217782649,3.5) coordinate (Spl2se1pa1) at (-7.33012701892,2.5) coordinate (Spl2se1pa2) at (-7.59807621135,1.5) coordinate (Spl2se1pa3) at (-7.86602540378,0.5) ;
\path coordinate (Spl3se1pa0) at (-7.86602540378,-0.5) coordinate (Spl3se1pa1) at (-7.59807621135,-1.5) coordinate (Spl3se1pa2) at (-7.33012701892,-2.5) coordinate (Spl3se1pa3) at (-7.06217782649,-3.5) ;
\path coordinate (Spl4se1pa0) at (-6.56217782649,-4.36602540378) coordinate (Spl4se1pa1) at (-5.83012701892,-5.09807621135) coordinate (Spl4se1pa2) at (-5.09807621135,-5.83012701892) coordinate (Spl4se1pa3) at (-4.36602540378,-6.56217782649) ;
\path coordinate (Spl1se2pa0) at (-3.5,-7.06217782649) coordinate (Spl1se2pa1) at (-2.5,-7.33012701892) coordinate (Spl1se2pa2) at (-1.5,-7.59807621135) coordinate (Spl1se2pa3) at (-0.5,-7.86602540378) ;
\path coordinate (Spl2se2pa0) at (0.5,-7.86602540378) coordinate (Spl2se2pa1) at (1.5,-7.59807621135) coordinate (Spl2se2pa2) at (2.5,-7.33012701892) coordinate (Spl2se2pa3) at (3.5,-7.06217782649) ;
\path coordinate (Spl3se2pa0) at (4.36602540378,-6.56217782649) coordinate (Spl3se2pa1) at (5.09807621135,-5.83012701892) coordinate (Spl3se2pa2) at (5.83012701892,-5.09807621135) coordinate (Spl3se2pa3) at (6.56217782649,-4.36602540378) ;
\path coordinate (Spl4se2pa0) at (7.06217782649,-3.5) coordinate (Spl4se2pa1) at (7.33012701892,-2.5) coordinate (Spl4se2pa2) at (7.59807621135,-1.5) coordinate (Spl4se2pa3) at (7.86602540378,-0.5) ;

  \fill[white,opacity=0.75] (-9,-9) rectangle (9,9);
  \foreach \player in {1,2,3,4} {
    \foreach \pa/\sh in {0/diamond, 1/star, 2/triang, 3/starplus} {
      \foreach \se in {0,1,2} {
        \pawn{\player}{\sh}{Spl\player se\se pa\pa}
      }
    }
  }
  \end{scope}
  \draw (-2,-2) rectangle (2,2);
\end{tikzpicture}
\caption{Pawn arrangment}
\label{fig:pawn-arrangment}
\end{figure}

\newpage

\noindent
\begin{tikzpicture}[remember picture, overlay,shift={(current page.center)}]
  \begin{scope}[yshift=9.5cm]
    %This file has been generated by POCcostFunction.py

%Places for cards
\matrix[column sep=5mm]{
\node[card place,draw=player1] {Player 1}; & \node[card place,draw=player2] {Player 2}; & \node[card place,draw=player3] {Player 3}; & \node[card place,draw=player4] {Player 4};\\
};

  \end{scope}
  \begin{scope}[yshift=-2.5cm]
    %This file has been generated by POCcostFunction.py

%Player sector colors
\path[fill=player1](8.8,0.0) -- (7.6210235533,4.4) -- (6.92820323028,4.0) -- (8.0,0.0) -- cycle;
\path[fill=player1](-4.4,7.6210235533) -- (-7.6210235533,4.4) -- (-6.92820323028,4.0) -- (-4.0,6.92820323028) -- cycle;
\path[fill=player1](-4.4,-7.6210235533) -- (-1.61653377487e-15,-8.8) -- (-1.46957615898e-15,-8.0) -- (-4.0,-6.92820323028) -- cycle;
\path[fill=player2](7.6210235533,4.4) -- (4.4,7.6210235533) -- (4.0,6.92820323028) -- (6.92820323028,4.0) -- cycle;
\path[fill=player2](-7.6210235533,4.4) -- (-8.8,1.07768918325e-15) -- (-8.0,9.79717439318e-16) -- (-6.92820323028,4.0) -- cycle;
\path[fill=player2](-1.61653377487e-15,-8.8) -- (4.4,-7.6210235533) -- (4.0,-6.92820323028) -- (-1.46957615898e-15,-8.0) -- cycle;
\path[fill=player3](4.4,7.6210235533) -- (5.38844591625e-16,8.8) -- (4.89858719659e-16,8.0) -- (4.0,6.92820323028) -- cycle;
\path[fill=player3](-8.8,1.07768918325e-15) -- (-7.6210235533,-4.4) -- (-6.92820323028,-4.0) -- (-8.0,9.79717439318e-16) -- cycle;
\path[fill=player3](4.4,-7.6210235533) -- (7.6210235533,-4.4) -- (6.92820323028,-4.0) -- (4.0,-6.92820323028) -- cycle;
\path[fill=player4](5.38844591625e-16,8.8) -- (-4.4,7.6210235533) -- (-4.0,6.92820323028) -- (4.89858719659e-16,8.0) -- cycle;
\path[fill=player4](-7.6210235533,-4.4) -- (-4.4,-7.6210235533) -- (-4.0,-6.92820323028) -- (-6.92820323028,-4.0) -- cycle;
\path[fill=player4](7.6210235533,-4.4) -- (8.8,0.0) -- (8.0,0.0) -- (6.92820323028,-4.0) -- cycle;
%Polygon
\draw[polygon] (8.0,0.0) -- (6.92820323028,4.0) -- (4.0,6.92820323028) -- (4.89858719659e-16,8.0) -- (-4.0,6.92820323028) -- (-6.92820323028,4.0) -- (-8.0,9.79717439318e-16) -- (-6.92820323028,-4.0) -- (-4.0,-6.92820323028) -- (-1.46957615898e-15,-8.0) -- (4.0,-6.92820323028) -- (6.92820323028,-4.0) -- cycle;
%Nodes and links
\path[start] (7.86602540378,0.5) coordinate[start] (s0) (7.59807621135,1.5) coordinate[start] (s1) (7.33012701892,2.5) coordinate[start] (s2) (7.06217782649,3.5) coordinate[start] (s3) (6.56217782649,4.36602540378) coordinate[start] (s4) (5.83012701892,5.09807621135) coordinate[start] (s5) (5.09807621135,5.83012701892) coordinate[start] (s6) (4.36602540378,6.56217782649) coordinate[start] (s7) (3.5,7.06217782649) coordinate[start] (s8) (2.5,7.33012701892) coordinate[start] (s9) (1.5,7.59807621135) coordinate[start] (s10) (0.5,7.86602540378) coordinate[start] (s11) (-0.5,7.86602540378) coordinate[start] (s12) (-1.5,7.59807621135) coordinate[start] (s13) (-2.5,7.33012701892) coordinate[start] (s14) (-3.5,7.06217782649) coordinate[start] (s15) (-4.36602540378,6.56217782649) coordinate[start] (s16) (-5.09807621135,5.83012701892) coordinate[start] (s17) (-5.83012701892,5.09807621135) coordinate[start] (s18) (-6.56217782649,4.36602540378) coordinate[start] (s19) (-7.06217782649,3.5) coordinate[start] (s20) (-7.33012701892,2.5) coordinate[start] (s21) (-7.59807621135,1.5) coordinate[start] (s22) (-7.86602540378,0.5) coordinate[start] (s23) (-7.86602540378,-0.5) coordinate[start] (s24) (-7.59807621135,-1.5) coordinate[start] (s25) (-7.33012701892,-2.5) coordinate[start] (s26) (-7.06217782649,-3.5) coordinate[start] (s27) (-6.56217782649,-4.36602540378) coordinate[start] (s28) (-5.83012701892,-5.09807621135) coordinate[start] (s29) (-5.09807621135,-5.83012701892) coordinate[start] (s30) (-4.36602540378,-6.56217782649) coordinate[start] (s31) (-3.5,-7.06217782649) coordinate[start] (s32) (-2.5,-7.33012701892) coordinate[start] (s33) (-1.5,-7.59807621135) coordinate[start] (s34) (-0.5,-7.86602540378) coordinate[start] (s35) (0.5,-7.86602540378) coordinate[start] (s36) (1.5,-7.59807621135) coordinate[start] (s37) (2.5,-7.33012701892) coordinate[start] (s38) (3.5,-7.06217782649) coordinate[start] (s39) (4.36602540378,-6.56217782649) coordinate[start] (s40) (5.09807621135,-5.83012701892) coordinate[start] (s41) (5.83012701892,-5.09807621135) coordinate[start] (s42) (6.56217782649,-4.36602540378) coordinate[start] (s43) (7.06217782649,-3.5) coordinate[start] (s44) (7.33012701892,-2.5) coordinate[start] (s45) (7.59807621135,-1.5) coordinate[start] (s46) (7.86602540378,-0.5) coordinate[start] (s47);
\path (-5.66458950116,-2.85371291234) coordinate[netnode] (n0) (-3.13953676463,-5.52085651601) coordinate[netnode] (n1) (2.28766112508,0.678831079224) coordinate[netnode] (n2) (5.70027506402,1.67571087853) coordinate[netnode] (n3) (0.162584255511,-2.39218233072) coordinate[netnode] (n4) (0.453712276129,-6.3681749908) coordinate[netnode] (n5) (-1.34171407598,5.95105537458) coordinate[netnode] (n6) (2.0758861129,-3.595689675) coordinate[netnode] (n7) (-6.04119876802,-1.06253756873) coordinate[netnode] (n8) (-1.37359793813,-5.95783698463) coordinate[netnode] (n9) (-2.05612586019,-3.73415858111) coordinate[netnode] (n10) (6.3509586276,-0.0416470987682) coordinate[netnode] (n11) (-4.30143454416,-4.09863362309) coordinate[netnode] (n12) (-2.25569745535,3.90707185545) coordinate[netnode] (n13) (-6.27206504806,0.823498351809) coordinate[netnode] (n14) (5.30372947143,3.47871124844) coordinate[netnode] (n15) (2.42292817113,5.84345970869) coordinate[netnode] (n16) (-3.81735339875,-1.91593227543) coordinate[netnode] (n17) (-1.73178248606,-1.64173770211) coordinate[netnode] (n18) (0.588091421565,6.10312705637) coordinate[netnode] (n19) (2.17912942633,-5.70194175276) coordinate[netnode] (n20) (2.03091755798,3.92679719021) coordinate[netnode] (n21) (3.8254204051,-2.28856060577) coordinate[netnode] (n22) (3.94086224942,4.70049587844) coordinate[netnode] (n23) (-5.57947267348,2.54232416471) coordinate[netnode] (n24) (0.00591003467172,-0.0102629739174) coordinate[netnode] (n25) (-0.124381953155,4.14846517207) coordinate[netnode] (n26) (5.28809431936,-3.57713120649) coordinate[netnode] (n27) (3.84836348643,-4.7382245695) coordinate[netnode] (n28) (-1.1949128307,2.06968815649) coordinate[netnode] (n29) (-3.53043878079,2.18202531352) coordinate[netnode] (n30) (0.916983837086,2.17886126415) coordinate[netnode] (n31) (3.56797064983,2.34784811551) coordinate[netnode] (n32) (4.26191591277,-0.0865436271976) coordinate[netnode] (n33) (-2.34545730164,0.295339841832) coordinate[netnode] (n34) (-3.20785167527,5.55624196324) coordinate[netnode] (n35) (5.84638346349,-1.78947744498) coordinate[netnode] (n36) (1.99038326562,-1.3369584776) coordinate[netnode] (n37) (-4.41618543933,0.204688715345) coordinate[netnode] (n38) (0.0690820574776,-4.4571725381) coordinate[netnode] (n39) (-4.48285212832,4.13754150742) coordinate[netnode] (n40);
\draw[link] (n38)  edge  (n14) (n38)  edge  (n17) (n38)  edge  (n24) (n38)  edge  (n34) (n39)  edge  (n4) (n39)  edge  (n5) (n39)  edge  (n7) (n39)  edge  (n10) (n39)  edge  (n18) (n30)  edge  (n13) (n30)  edge  (n24) (n30)  edge  (n29) (n30)  edge  (n34) (n30)  edge  (n38) (n30)  edge  (n40) (n31)  edge  (n2) (n31)  edge  (n21) (n31)  edge  (n25) (n31)  edge  (n26) (n31)  edge  (n29) (n31)  edge  (n32) (n32)  edge  (n3) (n32)  edge  (n15) (n32)  edge  (n33) (n33)  edge  (n3) (n33)  edge  (n11) (n33)  edge  (n37) (n34)  edge  (n17) (n34)  edge  (n25) (n35)  edge  (s14) (n35)  edge  (s15) (n7)  edge  (n37) (s44)  edge  (n27) (s44)  edge  (n36) (s45)  edge  (n27) (s45)  edge  (n36) (s46)  edge  (n11) (s46)  edge  (n36) (s47)  edge  (n11) (s47)  edge  (n36) (n12)  edge  (s28) (n12)  edge  (s29) (n12)  edge  (s30) (n12)  edge  (s31) (n12)  edge  (n10) (n12)  edge  (n17) (n13)  edge  (n6) (n13)  edge  (n26) (n13)  edge  (n29) (n13)  edge  (n35) (n13)  edge  (n40) (n10)  edge  (n4) (n10)  edge  (n17) (n10)  edge  (n18) (n11)  edge  (n36) (s7)  edge  (n16) (s24)  edge  (n14) (s4)  edge  (n15) (n18)  edge  (n4) (n18)  edge  (n17) (n18)  edge  (n25) (n18)  edge  (n34) (n19)  edge  (n26) (s19)  edge  (n24) (s18)  edge  (n35) (s13)  edge  (n6) (s13)  edge  (n19) (s12)  edge  (n6) (s12)  edge  (n19) (s11)  edge  (n16) (s11)  edge  (n19) (s10)  edge  (n16) (s10)  edge  (n19) (s17)  edge  (n35) (s16)  edge  (n35) (s38)  edge  (n5) (s35)  edge  (n5) (s35)  edge  (n9) (s34)  edge  (n5) (s34)  edge  (n9) (s37)  edge  (n5) (s37)  edge  (n20) (s36)  edge  (n5) (s36)  edge  (n20) (s22)  edge  (n14) (s22)  edge  (n24) (s9)  edge  (n16) (s9)  edge  (n19) (s8)  edge  (n16) (s23)  edge  (n14) (s3)  edge  (n3) (s3)  edge  (n15) (s2)  edge  (n3) (s2)  edge  (n15) (s1)  edge  (n3) (s1)  edge  (n11) (s0)  edge  (n3) (s0)  edge  (n11) (s6)  edge  (n15) (s5)  edge  (n15) (n40)  edge  (s16) (n40)  edge  (s17) (n40)  edge  (s18) (n40)  edge  (s19) (n29)  edge  (n25) (n29)  edge  (n26) (n29)  edge  (n34) (n28)  edge  (s39) (n28)  edge  (s40) (n28)  edge  (s41) (n28)  edge  (s42) (n28)  edge  (s43) (n28)  edge  (n7) (n23)  edge  (s4) (n23)  edge  (s5) (n23)  edge  (s6) (n23)  edge  (s7) (n23)  edge  (s8) (n23)  edge  (n21) (n22)  edge  (n7) (n22)  edge  (n27) (n22)  edge  (n33) (n22)  edge  (n36) (n22)  edge  (n37) (n21)  edge  (n16) (n21)  edge  (n26) (n21)  edge  (n32) (n20)  edge  (s38) (n20)  edge  (s39) (n20)  edge  (s40) (n20)  edge  (n7) (n27)  edge  (s41) (n27)  edge  (s42) (n27)  edge  (s43) (n26)  edge  (n6) (n25)  edge  (n37) (n8)  edge  (s23) (n8)  edge  (s24) (n8)  edge  (s25) (n8)  edge  (s26) (n8)  edge  (s27) (n8)  edge  (n38) (n9)  edge  (s32) (n9)  edge  (s33) (n9)  edge  (n39) (s20)  edge  (n14) (s20)  edge  (n24) (s21)  edge  (n14) (s21)  edge  (n24) (n0)  edge  (s25) (n0)  edge  (s26) (n0)  edge  (s27) (n0)  edge  (s28) (n0)  edge  (s29) (n0)  edge  (n17) (n1)  edge  (s30) (n1)  edge  (s31) (n1)  edge  (s32) (n1)  edge  (s33) (n1)  edge  (n9) (n1)  edge  (n10) (n2)  edge  (n21) (n2)  edge  (n25) (n2)  edge  (n32) (n2)  edge  (n33) (n2)  edge  (n37) (n4)  edge  (n7) (n4)  edge  (n25) (n4)  edge  (n37) (n6)  edge  (s14) (s15)  edge  (n6);
%Named coordinates for starting locations
\path coordinate (Spl1se0pa0) at (7.86602540378,0.5) coordinate (Spl1se0pa1) at (7.59807621135,1.5) coordinate (Spl1se0pa2) at (7.33012701892,2.5) coordinate (Spl1se0pa3) at (7.06217782649,3.5) ;
\path coordinate (Spl2se0pa0) at (6.56217782649,4.36602540378) coordinate (Spl2se0pa1) at (5.83012701892,5.09807621135) coordinate (Spl2se0pa2) at (5.09807621135,5.83012701892) coordinate (Spl2se0pa3) at (4.36602540378,6.56217782649) ;
\path coordinate (Spl3se0pa0) at (3.5,7.06217782649) coordinate (Spl3se0pa1) at (2.5,7.33012701892) coordinate (Spl3se0pa2) at (1.5,7.59807621135) coordinate (Spl3se0pa3) at (0.5,7.86602540378) ;
\path coordinate (Spl4se0pa0) at (-0.5,7.86602540378) coordinate (Spl4se0pa1) at (-1.5,7.59807621135) coordinate (Spl4se0pa2) at (-2.5,7.33012701892) coordinate (Spl4se0pa3) at (-3.5,7.06217782649) ;
\path coordinate (Spl1se1pa0) at (-4.36602540378,6.56217782649) coordinate (Spl1se1pa1) at (-5.09807621135,5.83012701892) coordinate (Spl1se1pa2) at (-5.83012701892,5.09807621135) coordinate (Spl1se1pa3) at (-6.56217782649,4.36602540378) ;
\path coordinate (Spl2se1pa0) at (-7.06217782649,3.5) coordinate (Spl2se1pa1) at (-7.33012701892,2.5) coordinate (Spl2se1pa2) at (-7.59807621135,1.5) coordinate (Spl2se1pa3) at (-7.86602540378,0.5) ;
\path coordinate (Spl3se1pa0) at (-7.86602540378,-0.5) coordinate (Spl3se1pa1) at (-7.59807621135,-1.5) coordinate (Spl3se1pa2) at (-7.33012701892,-2.5) coordinate (Spl3se1pa3) at (-7.06217782649,-3.5) ;
\path coordinate (Spl4se1pa0) at (-6.56217782649,-4.36602540378) coordinate (Spl4se1pa1) at (-5.83012701892,-5.09807621135) coordinate (Spl4se1pa2) at (-5.09807621135,-5.83012701892) coordinate (Spl4se1pa3) at (-4.36602540378,-6.56217782649) ;
\path coordinate (Spl1se2pa0) at (-3.5,-7.06217782649) coordinate (Spl1se2pa1) at (-2.5,-7.33012701892) coordinate (Spl1se2pa2) at (-1.5,-7.59807621135) coordinate (Spl1se2pa3) at (-0.5,-7.86602540378) ;
\path coordinate (Spl2se2pa0) at (0.5,-7.86602540378) coordinate (Spl2se2pa1) at (1.5,-7.59807621135) coordinate (Spl2se2pa2) at (2.5,-7.33012701892) coordinate (Spl2se2pa3) at (3.5,-7.06217782649) ;
\path coordinate (Spl3se2pa0) at (4.36602540378,-6.56217782649) coordinate (Spl3se2pa1) at (5.09807621135,-5.83012701892) coordinate (Spl3se2pa2) at (5.83012701892,-5.09807621135) coordinate (Spl3se2pa3) at (6.56217782649,-4.36602540378) ;
\path coordinate (Spl4se2pa0) at (7.06217782649,-3.5) coordinate (Spl4se2pa1) at (7.33012701892,-2.5) coordinate (Spl4se2pa2) at (7.59807621135,-1.5) coordinate (Spl4se2pa3) at (7.86602540378,-0.5) ;

  \end{scope}
\end{tikzpicture}

\newpage
\section*{Pawns}
\begin{tikzpicture}
  \newcommand{\colstride}{4.5}
  \newcommand{\rowstride}{1.5}
  \foreach \p in {1,2,3,4} {
    \foreach \xx/\sh in {0/diamond, 1/star, 2/triang, 3/starplus} {
      \foreach \x in {0,1,2,3} {
        \pawn{\p}{\sh}{\colstride*\xx+\x,-\rowstride*\p}
      }
    }
  }
\end{tikzpicture}
\section*{Letter tokens}
\begin{tikzpicture}
  \newcommand{\colstridetoken}{2cm}
  \newcommand{\rowstridetoken}{2cm}
  \newcommand{\tokensize}{\cardwidth/2.2}
  \foreach \p in {1,2,3,4} {
    \foreach \xx/\sh in {0/diamond,1/star,2/triang,3/starplus} {
      \node[circle,fill=player\p,minimum size=\tokensize] at (\colstridetoken*\xx,-\rowstridetoken*\p) {};
      \node[inner sep=0pt,\sh,fill=white,minimum size=0.8*\tokensize] at (\colstridetoken*\xx,-\rowstridetoken*\p) {};
    }
  }
\end{tikzpicture}

\section*{Word cards}
See next pages.

\tikzset{ 
  table/.style={
    matrix of nodes,
    row sep=-\pgflinewidth,
    column sep=-\pgflinewidth,
    nodes={
      rectangle,
      draw=black,
      align=center
    },
    minimum height=\cardheight,
    minimum width=\cardwidth,
    text depth=0.5ex,
    text height=2ex,
    nodes in empty cells,
    font=\sffamily
  }
}

\newpage\noindent
\begin{tikzpicture}
\matrix[table]{
TEST & WITH & HILL & BIRD & CORN\\
BEAR & SURE & MOST & GRAY & FINE\\
FIND & KILL & HOLE & FORM & ABLE\\
PART & LAND & ROSE & TONE & FIRE\\
THIN & WENT & SING & FALL & SHOP\\
};
\end{tikzpicture}
\begin{tikzpicture}
\matrix[table]{
DUCK & JUMP & ONCE & ROOM & LADY\\
LONG & PLAN & WEST & SOFT & KEEP\\
HERE & TERM & LEAD & GONE & PULL\\
PATH & EVEN & INCH & STAY & BANK\\
KEPT & MISS & BASE & CROP & THEN\\
};
\end{tikzpicture}
\begin{tikzpicture}
\matrix[table]{
CITY & EASE & KNOW & FEEL & WEEK\\
WOOD & FELT & HUGE & FLAT & LAST\\
OVER & CAMP & LATE & PORT & SEND\\
NEAR & MUCH & FOOD & RAIN & IDEA\\
GIVE & SHOW & FEAR & BURN & SEED\\
};
\end{tikzpicture}


\section*{Acknowledgements}
The development of \gamename{} was made possible by the support of grant IIS-1717656 by the National Science Foundation.

\includegraphics[width=3cm]{figures/NSF-4-Color-bitmap-Logo}

\end{document}
